\pg{351} ente, que é \emph{acto puro}, não encerra em si mesmo nenhuma potencialidade. Ora o ente, que não encerra em si nenhuma potencialidade, é absolutamente \emph{simples}; porque é só acto, e o acto não tem partes, em que possa dividir-se. Logo o ente, que é \emph{acto puro}, é absolutamente \emph{simples}.

b) \emph{O ente, que é absolutamente simples, é acto puro.} --- O ente, que é absolutamente \emph{simples}, não consta de partes, e por isso não é composto de potencia e de acto. Ora o ente, que não é composto de potencia e de acto, é \emph{acto puro}. Logo o ente, que é absolutamente \emph{simples}, é \emph{acto puro} \footnote{O \emph{acto puro} é absolutamente \emph{simples}. O acto é perfeição, e o acto puro é só perfeição, é perfeição subsistente; e por isso exclue toda a potencialidade, não pode ser composto de potencia e de acto, é absolutamente simples. Por quanto, como diz S. Thomaz, em todo o composto deve encontrar-se o acto e a potencia. Com effeito, muitas coisas não podem tornar-se simplesmente uma só coisa, se não emquanto uma seja acto e outra potencia; visto que, se fossem todas em acto, haveria um aggregado de coisas, mas nunca uma só coisa composta. Ora o acto puro, excluindo toda e qualquer potencialidade, é absolutamente simples. --- Não é menos evidente que o ente absolutamente \emph{simples} é \emph{acto puro}. Por quanto, o ente absolutamente \emph{simples}, que subsiste perfeitamente em si mesmo, ou é só potencia, ou é só acto. Ora é absurdo dizer que é só potencia; pois, como vimos, a potencia, denotando limitação e imperfeição, não pode subsistir em si mesma. Logo deve ser só acto, deve ser acto puro (\emph{Cf}. \emph{C}. \emph{Gent}., I, 18).}.

\textbf{75. O ente composto não pode ser acto puro.} --- Ente \emph{composto} é o que tem partes. Ora o que tem partes não pode ser \emph{acto puro}; porque --- ou uma das partes será \emph{acto} em relação a outra, --- ou todas as partes serão como que \emph{em potencia} em relação ao todo. Logo o ente \emph{composto} não pode ser \emph{acto puro} \footnote{Diz S. Thomaz: «In omni composito oportet esse potentiam et actum; quia vel una pars est actus respectu alterius, vel saltem omnes partes sunt in potentia respectu totius» (\emph{Sum. Th.}, p. I, q. 3, a. 7).}.

\textbf{76. O ente, que é absolutamente simples, é necessario e infinito; o ente, que é necessario e infinito, é absolutamente simples.}

a) \emph{O ente, que é absolutamente simples, é necessario e infinito.} --- O ente, que é absolutamente \emph{simples}, é acto puro, é o proprio ser subsistente. Ora o acto puro, o proprio ser subsistente, é \emph{necessario} e \emph{infinito}. Logo o ente que é absolutamente \emph{simples}, é \emph{necessario} e \emph{infinito}.

b) \emph{O ente, que é necessario e infinito, é absolutamente simples.} --- O ente, que é \emph{necessario} e \emph{infinito}, possue a existencia